
En
effet, si \(r=0\), \(L\) est évidemment un \(a_0\)-faisceau. Si \(r\ge1\), par 2.4, \(h^0(L(-1))=0\) puisque \(d(L(-1))<0\). Soit \(Y\) une section hyperplane de \(X\) et soit \(0\longrightarrow L(-1)\longrightarrow L\longrightarrow L_Y\longrightarrow0\) une suite exacte définie par \(Y\). Si \(r=1\), par 1.7, \(h^0(L_Y(-1))=a_1\); donc \(L\) est un \((0,a_1)\)-faisceau. Si \(r\ge2\), on peut prendre \(Y\) intègre grâce au théorème de Bertini et on peut ainsi supposer par récurrence que \(L_Y\) est un \((0,\ldots,a_r)\)-faisceau; d'où l'assertion.

\textbf{Lemme 2.7.} \emph{Si \(r=2\) et \(d(L)\le0\), alors} :
\begin{enumerate}[label=(\roman*)]
\item \emph{Pour tout \(n\), \(h^1(L(n))\le a_1(a_1+1)a_2-a_0\)}
\item \emph{\(a_0\le a_1(a_1+1)\).}
\end{enumerate}

Soit \(Y\) une section hyperplane intègre de \(X\). Par 1.7 et 2.6, \(L_Y\) est un \((0,a_2)\)-faisceau à polynôme de Hilbert \(\chi(L_Y(n))=a_1+a_2\binom{n+1}{1}\). Donc par 1.11, \(-a_1\ge0\) et \(L_Y\) est \([-a_1]\)-régulier.

Considérons la suite exacte suivante :
\[
H^0(L_Y(n))\longrightarrow H^1(L(n-1))
\longrightarrow H^1(L(n))\longrightarrow H^1(L_Y(n)).
\]
Pour \(n\ge -a_1-1\), on a \(h^1(L(-a_1-2))\ge h^1(L(n))\) car \(h^1(L_Y(n))=0\) par 1.3 (i); pour \(n\le -1\), on a \(h^1(L(n-1))\le h^1(L(-1))\) car \(h^0(L_Y(n))=0\) par 2.4. Enfin, pour tout \(n\), \(h^0(L_Y(n-1))\le h^0(L_Y(n))\) car il existe une injection \(L_Y(n-1)\longrightarrow L_Y(n)\), et \(h^0(L_Y(-a_1-1))=a_2\binom{-a_1}{1}+a_1\) car \(h^1(L_Y(-a_1-1))=0\); donc pour \(n\le -a_1-2\), on a \(h^1(L(n))-h^1(L(n+1))\le -a_2a_1+a_1\). Comme conséquence, on en conclut que pour tout \(n\),
\[
h^1(L(n))\le h^1(L(-a_1-2))+(-a_1-1)(-a_2a_1+a_1).
\tag{2.7.1}
\]

